\documentclass{ctexart}
\usepackage{graphicx} % Required for inserting images
\usepackage[margin=2.5cm]{geometry}
\usepackage{booktabs}
\usepackage{float}
\usepackage{soulutf8}
\usepackage{enumitem}
\usepackage{pifont}
\usepackage{pgfplots}
\usepackage{longtable}
\pgfplotsset{compat=newest}
\usepackage{tikz}
\usepackage{xcolor} % 用于颜色设置
\usepackage{listings} % 用于代码块
\usepackage{fontspec} % 关键：用于使用系统字体 (XeLaTeX/LuaLaTeX)
% Linux 环境中没有 Windows 宋体时，使用可用的中文衬线字体作为宋体替代
\setCJKmainfont{Noto Serif CJK SC}
\usepackage{authblk}
\usepackage[
    colorlinks=true,      % 使用颜色而非方框
    linkcolor=blue,       % 内部链接颜色
    citecolor=green,      % 引用链接颜色
    urlcolor=red,         % URL 链接颜色
    bookmarks=true,       % 生成 PDF 书签
    bookmarksopen=true,   % 默认展开书签
    pdfstartview=FitH     % 页面宽度适应窗口
]{hyperref}

\sethlcolor{yellow}%设置所需高亮颜色
\usepackage[most]{tcolorbox}
\tcbuselibrary{listings,skins,breakable}
% 页眉页脚
\usepackage{fancyhdr}
\pagestyle{fancy}

% 定义草绿色
\definecolor{grassgreen}{RGB}{124,179,66}
% 补全缺失的浅蓝色定义（防止编译报错）
\definecolor{lightblue}{RGB}{30,144,255} 

% 清空默认
\fancyhf{}
\setlength{\headheight}{14pt}

% 页眉
\fancyhead[L]{\textcolor{lightblue}{\kaishu Computer Vision}}
\fancyhead[R]{\textcolor{lightblue}{\textit{2026.9-2027.1}}}

% 页脚页码居中
\fancyfoot[C]{\textcolor{lightblue}{\thepage}}

% 页眉、页脚线颜色
\renewcommand{\headrulewidth}{1pt}
\renewcommand{\footrulewidth}{1pt}
\renewcommand{\headrule}{\color{blue!30!white}\hrule height 1pt width\headwidth}
\renewcommand{\footrule}{\color{blue!30!white}\hrule height 1pt width\headwidth}
\newtcolorbox{fancybox}[2][blue]{ % 可选颜色参数
  enhanced,
  breakable,
    colback=#1!3!white,
    colframe=#1!18!white,
    boxrule=0.8pt,
    arc=6pt,
    outer arc=6pt,
  left=10pt,right=10pt,top=12pt,bottom=10pt,
  title={\textbf{#2}},
    fonttitle=\bfseries\color{#1!65!black},
  attach boxed title to top left={yshift=-2mm,xshift=5mm},
  boxed title style={
        colback=#1!8!white,
        colframe=#1!18!white,
        boxrule=0.6pt,
        arc=4pt,
  },
    shadow={1.5pt}{-1.5pt}{0pt}{black!10},
}

% 定义一套适合C++的配色方案
\usepackage{tabularx}
\newfontfamily{\consolasfont}{Ubuntu}

\newcommand{\bluecode}[1]{{\color{blue}\consolasfont #1}}
\lstdefinestyle{MyCppStyle}{
    language = C++,
    % 确保这里使用 \consolasfont
    basicstyle = \consolasfont\small, 
    keywordstyle = \color{blue!70!black}, 
    commentstyle = \color{green!50!black}, 
    stringstyle = \color{red!60!black}, 
    numberstyle = \tiny\color{gray}, 
    numbers = left, 
    frame = single, 
    rulecolor = \color{lightgray}, 
    breaklines = true, 
    backgroundcolor = \color{white}, 
    tabsize = 4, 
    showstringspaces = false, 
    xleftmargin = 2em, 
    escapeinside=``,
}
% --- lstdefinestyle HTML 风格定义 ---
\lstdefinestyle{mhtml}{
    language = HTML,
    basicstyle = \consolasfont\small, 
    tagstyle = \color{blue!70!black}\bfseries,
    keywordstyle = \color{purple!80!black},
    commentstyle = \color{green!50!black}\itshape, 
    stringstyle = \color{red!60!black}, 
    numberstyle = \tiny\color{gray}, 
    numbers = left, 
    frame = single, 
    rulecolor = \color{lightgray}, 
    breaklines = true, 
    backgroundcolor = \color{white}, 
    tabsize = 4, 
    showstringspaces = false, 
    xleftmargin = 2em, 
    escapeinside=``,
    morekeywords={als, doctype, html, head, title, meta, link, body, div, span, h1, h2, h3, h4, h5, h6, p, a, img, ul, ol, li, table, tr, th, td, pre, code, script, style},
    morekeywords=[2]{class, id, href, src, rel, type, charset, lang, name, content, style},
    keywordstyle=[2]\color{purple!80!black},
}
\usepackage{hyperref}
\hypersetup{
    colorlinks=true,      
    linkcolor=blue,       
    filecolor=magenta,    
    urlcolor=cyan,        
    pdftitle={我的文档标题}, 
}

% ==================== 正文部分 ====================
\begin{document}

% 开启无页眉页脚模式（针对封面和目录）
\pagestyle{empty}

% ------------------ 美化封面页 ------------------
\begin{titlepage}
    \centering
    % 顶部 TikZ 几何几何装饰条（全宽 bleed 效果）
    \begin{tikzpicture}[remember picture, overlay]
        \fill[blue!50!white] (current page.north west) rectangle ([yshift=-2cm]current page.north east);
        \fill[grassgreen] (current page.north west) ++(0,-2cm) rectangle ([yshift=-2.2cm]current page.north east);
    \end{tikzpicture}
    
    \vspace*{3.5cm}
    
    % 大标题
    {\fontsize{30pt}{36pt}\selectfont\bfseries\color{blue!80!black} 计算机视觉（CV）}\\[0.6em]
    % 副标题或英文标识
    {\fontsize{14pt}{18pt}\selectfont\color{gray}\textsf{Computer Vision}}\\[1cm]
    
    % 居中装饰线
    \centering\makebox[4cm]{\color{grassgreen}\rule{5cm}{1.5pt}}
    
    \vfill
    
    % 使用 tcolorbox 包装的精致作者信息框
    \begin{tcolorbox}[
        enhanced,
        width=0.55\textwidth,
        colback=blue!5!white,
        colframe=blue!5!white,
        boxrule=1.2pt,
        arc=6pt,
        sharp corners=downhill, % 别致的对角切角效果
        shadow={2pt}{-2pt}{0pt}{black!20},
        halign=center,
        valign=center,
        top=15pt, bottom=15pt
    ]
        {\large\bfseries\color{black!80} Author：} {\large\kaishu Simson} \\[0.6em]
        {\large\bfseries\color{black!80} Date：} {\large\kaishu September 2026}
    \end{tcolorbox}
    
    \vfill
    \vspace*{2cm}
    
    % 底部标识
    {\color{gray!80}\small\the\year\ \copyright\ All Rights Reserved.}
\end{titlepage}

\pagestyle{fancy}

\newpage
\tableofcontents
\newpage

% 设置从当前页（正文第一页）开始采用阿拉伯数字编码，并自动重置为第 1 页
\pagenumbering{arabic}

\section{导论}

\textbf{可视计算的定义：}针对视觉内容生成与处理的所有计算机课学分制的总称（计算机图形学，计算机视觉，可视化，AR,VR）人类智能在感知，认知，决策，行动，感知的循环中逐渐积累。

重要课题：
\begin{itemize}
    \item 物理世界仿真与模拟
    \item 计算机图形学
    \item 计算机视觉
\end{itemize}

\section{VCL基础知识和技能}

\subsection{颜色，颜色感知和可视化}

人感知颜色靠的是对不同波长的光敏感的三种视锥细胞，这提供了颜色离散表达的基础。三种感光开展出RGB三色系统。颜色的匹配函数如下，是一个积分的办法。

\[
R=\int_0^{\infty}I(\lambda)\bar{r}(\lambda)d\lambda
\]
\[
G=\int_0^{\infty}I(\lambda)\bar{g}(\lambda)d\lambda
\]
\[
B=\int_0^{\infty}I(\lambda)\bar{b}(\lambda)d\lambda
\]

三种通道和原色的光由离散合成有两种机制：基于发光叠加相加混光（显示器发光原理）和基于白光吸收的相减混光（涂料和打印原理）。

计算机视觉基于人对光的感知。在人的感知之下，红色有前进的效应，蓝色会有看起来后退的效应。此即所谓色彩立体视觉。同时在人的认知下颜色也有冷色和暖色之分（这真的是课堂PPT里面的）。颜色信息在人眼中具有恒常性，不论在何种光照之下人对于颜色的认知不会发生大的改变。

当然这种逻辑感知也会带来相应的视觉错觉，譬如色诱导，使用补色可以让人自动脑补对应的颜色；再比如说边缘效应和移动错觉（动态立体视觉）。

人眼对于不同的光敏感性不同，对不同波段的光，施加灰度和模糊处理可以得到不同的认知效果（灰度指的是介于黑白之间的明暗程度）。

\section{显示（Display）}

\subsection{显示器的种类和发展历史}

\begin{enumerate}
    \item 矢量显示（阴极射线管）
    \item 二维显示（电子束）
    \begin{itemize}
        \item 像素：每个可以点亮的屏幕点
        \item 分辨率：品目像素格式
        \item 纵横比：像素列数除以行数
        \item 帧：显示的一个画面
        \item 隔行扫描：保证亮度均衡
    \end{itemize}
    \item 彩色二维显示（彩色荧光幕，不同荧光剂分区涂布）
    \item 二维显示（LCD）：统一发光，用液晶部分控制光的偏振，调节亮度
    \item 二维显示（OLED）：有机发光半导体显示器。独立的LED单元。
\end{enumerate}

\subsection{色域}

美国国家电视标准委员会的色域标准称为NTSC。

\begin{itemize}
    \item sRGB：微软，HP（72% NSTC）
    \item ARGB: Adobe颜色（设计软件）
    \item Rec.2020:4K(电视的标准)
    \item DCI-P3 : Apple标准
\end{itemize}

\subsection{像素标准}

显卡的存储和刷新越厉害，达到4K8K就越容易。

\subsection{立体显示}

立体显示依据的原理是双目视觉。所有的立体显示技术基于人类的生理特征。最开始的显示技术是双目偏振片的3D显示。沉浸式的3D显示技术（AR,VR,XR）跟踪头显，裸眼3D+眼睛跟踪。现有的技术向MR-mix方向进行发展。

裸眼3D技术：LCD屏幕的方法：微柱透镜显示技术。利用视差屏障，令OLED屏幕上不同的像素点用不同角度进入人眼。

全息显示技术（Hplographic display）：用干涉记录波前的信息（数字全息摄影），用衍射记录物光波场（受限于较小的范围）。

多模态显示技术：触觉显示（触觉输出），嗅觉显示（控制挥发的成分。。。？）

\section{图像（Image Representation）}

在二维空间中，每个点可以由四个参数表示。$E(x,y,\lambda,t)$其中$\lambda$字母的含义为光的波长。从数学上来说图像本身是一个连续的表征（真实世界），在屏幕上离散化显示。VCL的任务在于用离散方式进行还原。

\subsection{图像的表示}

图像显示的两个方法是矢量表示法（Vector Representation，对应SVG格式）和光栅化表示法（Raster Representation，像素级存储）。

矢量表示法的精确度高，用较小的空间，不会受到分辨率影响（可以用新的矢量计算适应不同屏幕不同大小）多用于图标等规则图形。Raster存储空间较大，受限于分辨率，但是可以表达更加不规则的图形。

\subsection{图形处理器（GPU）}

GPU有成百上千个简单的计算单元，可以在同一时间并行运行大量程序，比如计算每个像素的颜色，而CPU一般只有个位数的可并行的计算单元。故而GPU显卡适合处理进行大量计算的图形编程任务。

Pixel存储中有重要概念帧缓冲区（Frambuffers for image processing and display），指的是储存二维像素矩阵的内存大小。硬件以60Hz速度扫描帧缓冲区中的内容，这也决定了显示器的种类。

同样存储彩色图片的时候，Fullcolor（RGB），决定了显示器的颜色存储位数。一般来说真彩色图像的显示屏为24-bit位存储。如果颜色粗出空间较小，色彩的丰富度就会下降。

gif颜色存储的性质：Colormap映射。对于每一个像素点，存储一个8-bit常数，映射到RGB颜色表格Colormap，根据映射进行显示。在预先选好的表格中进行查找和压缩。

\begin{fancybox}{几种图片的文件格式}

\textbf{常见图像文件格式的全称：}
\begin{itemize}[leftmargin=1.6em,itemsep=1pt,topsep=3pt,parsep=0pt]
    \item \textbf{BMP}：Full color image（全彩色图像）
    \item \textbf{JPEG}：Joint Photographic Experts Group Format
    \item \textbf{TIFF}：Tagged-Image File Format
    \item \textbf{GIF}：CompuServe Graphics Interchange Format
    \item \textbf{PPM}：Portable PixMap Format（ASCII 或二进制）
    \item \textbf{EPS}：Encapsulated PostScript Format（ASCII）
\end{itemize}

\vspace{2pt}
\begin{center}
\small
\begin{tabular}{lccc}
\toprule
 & \textbf{BITS PER PIXEL} & \textbf{FILE SIZE} & \textbf{COMMENTS} \\
\midrule
JPEG & 24            & small  & lossy compression \\
TIFF & 8, 24         & medium & good general purpose \\
GIF  & 1, 4, 8       & medium & popular, but 8-bit \\
PPM  & 24            & big    & easy to read/write \\
EPS  & 1, 2, 4, 8, 24 & huge  & good for printing \\
\bottomrule
\end{tabular}
\end{center}
\end{fancybox}

同时，相较于24-bit的图像，有的高级图像会膨胀到每个像素96-bit，有一些其他的通道来表示图像。

\begin{enumerate}
    \item Alpha通道（8-bit）：透明度（Transparency），存储透明度常数，存在前中后背景。
    
    根据参数$\alpha$计算像素点的表示：$b = (1-\alpha) a_1 + \alpha a_2$
    \item Z-buffer（16-bit）：保存像素点所对应的3D空间点的深度。判断在叠加一层的时候应该占据多少通道和Alpha信息。
    
    \item 双通道表示：既要保存图像的现有数据，又要保存图像的旧信息，方便修改。
\end{enumerate}

\subsection{图像处理（Image Processing）}

图像处理按操作的作用范围可以分为两类基本操作。\textbf{点处理（Point Processing）}逐像素地修改像素值，输出仅是该像素自身像素值的函数；\textbf{滤波（Filtering）}的输出则是该像素周围（通常是局部）邻域（local neighborhood）内像素的函数，因而需要用到邻域中其他像素的信息。
我们也可以理解为一个是逐点操作，一个是区域化的操作（图像抖动，图像过滤，区域补全）。

除这两类基本操作之外，图像处理还包括以下其他课题：

\begin{itemize}
    \item 图像变换（Image transformation）：缩放（resize）、形变（warp）
    \item 图像压缩（Image compression）：无损（lossless）与有损（lossy）
    \item 图像编解码（Image encoding and decoding）：用于高效传输（efficient transmission）
    \item 图像分割（Image segmentation）
    \item 图像抠图（Image matting）
    \item 图像理解（Image understanding）：进一步走向计算机视觉（computer vision）
    \item 视频处理（Video processing）：推广到图像序列（a sequence of images）
    \item ……
\end{itemize}

\end{document}